\documentclass{resume} % Keep resume.cls from the original Overleaf template
\usepackage[dvipsnames]{xcolor}
\usepackage{hyperref}
\usepackage[left=0.75in,top=0.55in,right=0.75in,bottom=0.35in]{geometry}

\name{Kairui Feng}
\address{Phone: +86 185 9965 1229 \quad Email: \href{mailto:3405402279@qq.com}{3405402279@qq.com}}

\definecolor{CarnegieMellonRed}{RGB}{196,18,48}
\renewenvironment{rSection}[1]{
\sectionskip
\textcolor{CarnegieMellonRed}{\MakeUppercase{#1}}
\sectionlineskip
\hrule
\begin{list}{}{\setlength{\leftmargin}{1.5em}}
\item[]
}{\end{list}}

\begin{document}

\begin{rSection}{Education}
\textbf{Zhejiang University} \hfill \textit{Expected Jun. 2027}\\
B.Eng. in Information Security, College of Computer Science and Technology\\
GPA: 3.96/4.30 \quad Weighted Average: 87.33/100 \quad IELTS Academic: 6.5
\end{rSection}

\begin{rSection}{Security-Related Core Coursework}
\textbf{Security Coursework Average: 92.1/100 (18.5 credits)}\\
\small Software Security Principles (99); Software Security Theory \& Practice (92); Network Security Theory \& Practice (92); AI Security (89); AI Security \& Ethics (89); Advanced Cryptology (92); System Security (97); Wireless/IoT Security (84); Privacy, Security \& Digital Currency (95).\normalsize
\end{rSection}

\begin{rSection}{Selected Systems Projects and Research Experience}
\textbf{Linux Kernel Security Analysis and Exploitation} \hfill \textit{Independent Security Practice}\\
\begin{itemize}
    \item Analyzed kernel coloring techniques through source-code reading.
\end{itemize}

\textbf{gdb-helper} \hfill \textit{Open-Source Debugging Tool}\\
\begin{itemize}
    \item Extended an open-source GDB plugin with DeepSeek API integration to translate natural-language CLI requests into multi-step native GDB workflows. \url{https://github.com/fkrcarry/gdb-helper}
\end{itemize}

\textbf{RISC-V Operating System on a Self-Designed CPU} \hfill \textit{ZJU Computer Systems Integrated Curriculum}\\
\begin{itemize}
    \item Independently built a RISC-V OS running on a self-designed CPU. \textbf{SYS1:} designed and validated a single-cycle RISC-V CPU in HDL.
    \item \textbf{SYS2:} implemented a basic kernel with interrupts and context switching; built an RV64I CPU with a five-stage pipeline, branch prediction, bus interface, forwarding, L1 cache, and privileged instructions.
    \item \textbf{SYS3 (Labs 1--5 and xpart MMU):} implemented boot and timer interrupts, preemptive kernel threads, virtual memory, user mode and system calls, page faults, process creation, and MMU functionality.
    \item Received the \textit{Outstanding Project Award} from Prof. Rui Chang and Prof. Wenbo Shen (9 recipients among approximately 150 students).
\end{itemize}

\textbf{C++ RISC-V Compiler} \hfill \textit{Independent Systems Project}\\
\begin{itemize}
    \item Developed a C++ compiler targeting RISC-V.
\end{itemize}

\textbf{Mobile GUI-Agent-Based iOS App Behavior-Transformation Detection System} \hfill \textit{Under Prof. Haitao Xu}\\
\begin{itemize}
    \item Developed a mobile GUI-agent-based detector for iOS apps that change behavior or presentation during interaction.
\end{itemize}

\textbf{MCP Server Security Vulnerability Discovery} \hfill \textit{Independent Security Research}\\
\begin{itemize}
    \item Identified and publicly disclosed an SSRF vulnerability in \texttt{imprvhub/}\allowbreak\texttt{mcp-browser-agent} (up to 0.8.0) through user-controlled fields in \texttt{CallToolRequestSchema}; CVE-2026-5607: \url{https://vuldb.com/source_cve/355398}
\end{itemize}

\textbf{Other Security Research} \hfill \textit{Zhejiang University}\\
\begin{itemize}
    \item Code LLM red-teaming under Prof. Kui Ren and Prof. Zhan Qin (patent pending); SEED Labs security experimentation under Prof. Wenliang Du.
\end{itemize}
\end{rSection}

\begin{rSection}{Teaching Experience}
\textbf{Teaching Assistant}, Zhejiang University: Computer Systems I (SYS1), Computer Systems II (SYS2), and Principles and Practice of Software Security.\\
Delivered theory and laboratory instruction; released and checked labs; graded coursework and quizzes.
\end{rSection}

\begin{rSection}{Honors and Awards}
\begin{itemize}
    \item University-Level Student Pacesetter Award, Zhejiang University (twice)
    \item Second Prize, National College Student Information Security Competition
    \item First Prize, Zhejiang Provincial College Student Information Security Competition
    \item Third Place, Pwn Track, ZJUCTF 2025 (Zhejiang University)
    \item Ranked No. 4 in the Pwn track on Zhejiang University's CTF platform
    \item Outstanding Project Award, Computer Systems Integrated Curriculum
\end{itemize}
\end{rSection}

\begin{rSection}{Technical Skills}
\begin{tabular}{@{} >{\bfseries}l @{\hspace{1.5ex}} p{0.70\textwidth}}
Programming Languages & C, C++, Python, SystemVerilog \\
Systems \& Tools & Linux command line, Git, Docker, \LaTeX, QEMU kernel debugging, binary debugging, kernel vulnerability reproduction/analysis, fuzzing, Linux kernel modules/drivers, kernel privilege-escalation exploitation \\
\end{tabular}
\end{rSection}

\end{document}
